\documentclass[11pt,journal,onecolumn]{article}
\usepackage[utf8]{inputenc}
\usepackage[margin=1.0in]{geometry}
\usepackage{amsmath,amssymb,amsfonts}
\usepackage{booktabs}
\usepackage{graphicx}
\usepackage{cite}
\usepackage{url}
\usepackage{tabularx}
\usepackage{setspace}
\usepackage{lineno}
\usepackage{hyperref}
\hypersetup{colorlinks=true, linkcolor=blue, citecolor=blue, urlcolor=blue}

\doublespacing
\linenumbers

\title{\textbf{Multi-Factor Adaptive Traffic Signal Control Under Directional Demand Asymmetry and Dynamic Flow: Formulation, Stability Constraints, and Confirmatory Simulation Evaluation}}

\author{
    \textbf{Prince Tagadiya} \\
    Department of Civil and Environmental Engineering, Intelligent Transportation Systems Laboratory
}

\date{October 2026}

\begin{document}

\maketitle

\begin{abstract}
Traffic signal control at urban intersections faces persistent trade-offs between dynamic responsiveness to fluctuating vehicular queues and operational stability across varied demand regimes. Traditional fixed-time plans perform reliably under symmetric, steady-state flows but degrade severely under directional asymmetry or temporal surges. Conversely, unconstrained adaptive algorithms frequently exhibit excessive phase hunting, high-frequency switching, and loss of effective capacity due to yellow clearance lost time. This study presents a multi-factor adaptive traffic signal control formulation that integrates instantaneous halting queue length, upstream arrival pressure, backward finite-difference queue growth, accumulated waiting burden, and continuous red-age fairness. Operational stability is enforced through a constrained state-machine combining NEMA minimum and maximum green bounds, an absolute starvation override, and a calibrated switching hysteresis barrier ($H = 0.20$). The controller was evaluated in Eclipse SUMO via TraCI across seven demand scenarios and five independent evaluation seeds ($N = 5$ matched simulation realizations per cell, 140 total runs), benchmarking against original 90-second fixed-time control (C0), Webster-calibrated 36-second fixed-time control (C1), and basic queue-based adaptive control (C2). Evaluated strictly on a pre-registered Cohort A population, the proposed controller demonstrates its primary advantage under directional asymmetry, reducing mean waiting time by 61.6\% to 63.2\% relative to Webster-calibrated fixed-time timing ($p_{\text{adj}} < 0.04$ under Benjamini-Hochberg FDR control). Furthermore, across all seven evaluated scenarios, the proposed controller achieves a statistically significant 4.6\% to 20.3\% reduction in phase transitions compared to naive queue control ($p < 0.028$), successfully suppressing phase hunting without sacrificing throughput. Under symmetric steady-state demand, calibrated fixed-time control matches or slightly exceeds the adaptive controller's waiting-time performance, confirming the hypothesis that multi-factor adaptive control provides targeted benefits under asymmetric and dynamic regimes rather than uniform superiority across all traffic states.
\end{abstract}

\textbf{Keywords:} Adaptive traffic signal control, intelligent transportation systems, traffic signal optimization, queue dynamics, demand asymmetry, SUMO microscopic simulation, phase switching stability.

\section{Introduction}
Urban traffic congestion remains a critical global challenge, exacerbating economic inefficiency, vehicular emissions, and commuter delays \cite{papageorgiou2003review}. Signalized intersections represent the primary bottlenecks within urban arterial networks. Conventionally, signal timing plans rely on pre-timed or fixed-time schedules derived from historical flow observations, such as Webster's seminal cycle length formulation \cite{webster1958traffic}. While fixed-time control offers operational simplicity and predictable coordination, it cannot adapt dynamically to stochastic queue accumulation, sudden temporal surges, or pronounced directional demand asymmetries.

To address these limitations, adaptive traffic signal control (ATSC) systems have been extensively investigated. Classical systems such as SCOOT \cite{hunt1981scoot} and SCATS \cite{sims1980sydney} adjust cycle splits incrementally based on detector occupancy. More recent advances employ model-predictive control \cite{wongpiromsarn2012distributed}, max-pressure formulations \cite{varaiya2013max}, and reinforcement learning \cite{wei2018intellilight, wei2019survey}. However, existing adaptive methods frequently suffer from two operational deficiencies:
\begin{enumerate}
    \item \textbf{High-Frequency Phase Hunting}: Controllers driven purely by instantaneous queue comparisons often toggle signal phases erratically when competing demands are closely matched, squandering valuable green time on mandatory yellow clearance transitions ($Y = 3\,\text{s}$).
    \item \textbf{Evaluation against Straw-Man Baselines}: In the literature, proposed adaptive controllers are frequently benchmarked solely against uncalibrated, excessively long static cycle plans (e.g., 90–120 seconds) without comparison against properly calibrated, minimal-cycle fixed-time baselines or intermediate density controllers \cite{genders2016using}.
\end{enumerate}

This paper addresses these gaps by developing, mathematically formalizing, and rigorously auditing a multi-factor adaptive signal control architecture. The controller combines five distinct operational features: (1) halting queue length, (2) advance upstream arrival pressure, (3) temporal queue growth finite-difference, (4) accumulated waiting burden, and (5) continuous red-age fairness. Crucially, phase switching is governed by a dual-threshold mechanism combining absolute safety boundaries (NEMA minimum and maximum green, starvation override) with a dimensionless switching hysteresis barrier ($H = 0.20$).

\subsection{Research Contributions}
The specific contributions of this work are fourfold:
\begin{enumerate}
    \item \textbf{Multi-Factor State-Feedback Formulation}: We formulate a dimensionless composite priority metric that incorporates existing queues, incoming platoons, developing congestion, delay equity, and continuous starvation prevention.
    \item \textbf{Hysteresis-Stabilized Switching Law}: We integrate an explicit switching hysteresis margin that prevents high-frequency phase toggling, proving that phase transitions can be substantially reduced without increasing overall vehicular delay.
    \item \textbf{Rigorous Matched-Pairs Simulation Protocol}: We execute a pre-registered simulation study in Eclipse SUMO eliminating pseudoreplication by treating the macro simulation realization ($N = 5$ seeds across 7 scenarios) as the independent experimental unit, isolating a pure post-warmup evaluation population (Cohort A).
    \item \textbf{Empirical Hypothesis Evaluation without Overclaiming}: We evaluate whether the controller provides universal superiority or regime-specific advantages, confirming that its primary value emerges under directional asymmetry and phase stabilization rather than uniform delay minimization across symmetric steady-state regimes.
\end{enumerate}

\section{Related Work}
\subsection{Fixed-Time and Webster-Based Signal Control}
Fixed-time signal control traces its mathematical foundation to Webster's method \cite{webster1958traffic}, which minimizes average intersection delay by calculating an optimum cycle length $C_0 = (1.5L + 5)/(1 - Y)$ based on critical approach flow ratios and lost clearance time $L$. When traffic flows are steady and symmetric, Webster-calibrated timing plans provide high throughput and zero detector-induced estimation variance. However, static cycle plans fail to adjust when arrival distributions diverge from historical averages \cite{newell1989theory}.

\subsection{Actuated and Queue-Based Control}
Semi-actuated and fully-actuated controllers utilize inductive loops or stop-bar detectors to extend green phases via gap-out logic \cite{nema2005ts2}. Naive queue-length controllers compare stopping queues directly across conflicting approaches, switching whenever the opposing queue exceeds the active queue \cite{miller1963computer}. However, without temporal damping or hysteresis, naive queue controllers exhibit severe phase hunting under balanced demands, triggering frequent clearance intervals that degrade intersection capacity \cite{gershenson2005self}.

\subsection{Adaptive and Network-Wide Strategies}
Dynamic systems such as OPAC \cite{gartner1983opac} and RHODES \cite{mirchandani2001real} introduced rolling-horizon optimization using upstream flow prediction. Varaiya formulated the max-pressure control paradigm \cite{varaiya2013max}, mathematically establishing throughput optimality in store-and-forward queueing networks. While max-pressure algorithms guarantee network stability under infinite queues, they may cause unacceptable waiting times on minor approaches in the absence of explicit fairness and starvation constraints \cite{kouvelas2014maximum}.

\subsection{Learning-Based and Self-Organizing Approaches}
Recent investigations have applied deep reinforcement learning (DRL) to traffic signal control \cite{arel2010reinforcement, wei2018intellilight, wei2019survey}. Despite promising simulation demonstrations, DRL methods often lack safety guarantees, require complex state-action spaces, exhibit sensitivity to stochastic reward formulation, and present interpretability barriers for municipal traffic engineering deployment \cite{genders2016using}. A comparison of existing signal control methodologies against the proposed formulation is presented in Table~\ref{tab:lit_comparison}.

\begin{table}[htbp]
\centering
\small
\caption{Methodological Comparison of Traffic Signal Control Paradigms}
\label{tab:lit_comparison}
\begin{tabularx}{\textwidth}{lccccccX}
\toprule
\textbf{Control Paradigm} & \textbf{Adaptive?} & \textbf{Queue} & \textbf{Arrivals} & \textbf{Growth} & \textbf{Waiting} & \textbf{Fairness} & \textbf{Switching Dynamics} \\
\midrule
Fixed-Time (Webster) \cite{webster1958traffic} & No & No & No & No & No & Static & Cyclic pre-timed \\
Actuated Gap-Out \cite{nema2005ts2} & Partial & Yes & No & No & No & Max Green & Green extension / gap-out \\
Basic Queue / Density \cite{miller1963computer} & Yes & Yes & No & No & No & Starve & Naive sign comparison \\
Max-Pressure \cite{varaiya2013max} & Yes & Yes & No & No & No & No & Pressure difference \\
Self-Organizing \cite{cools2013self} & Yes & Yes & Yes & No & No & Dynamic & Rule-based thresholds \\
Deep RL (IntelliLight) \cite{wei2018intellilight} & Yes & Yes & Yes & Implicit & Yes & Reward & Neural network policy \\
\textbf{Proposed Controller (C3)} & \textbf{Yes} & \textbf{Yes} & \textbf{Yes} & \textbf{Yes} & \textbf{Yes} & \textbf{Yes} & \textbf{Hysteresis + NEMA Constraints} \\
\bottomrule
\end{tabularx}
\end{table}

\section{Problem Formulation}
\subsection{Intersection Geometry and Traffic Flow}
Consider an isolated four-arm orthogonal intersection $\mathcal{J}_0$ regulating North-South ($d=\text{NS}$, lanes $i \in \{N, S\}$) and East-West ($d=\text{EW}$, lanes $i \in \{E, W\}$) vehicular movements. Each approach has length $L_{\text{link}} = 250\,\text{m}$, speed limit $v_{\text{max}} = 13.89\,\text{m/s}$ ($50\,\text{km/h}$), and a single travel lane. The controller executes at uniform discrete time steps $\Delta t = 1.0\,\text{s}$.

\subsection{State Observables Extracted via Simulation Interface}
At each step $t$, the controller extracts simulated microscopic traffic observables via TraCI:
\begin{enumerate}
    \item \textbf{Halting Queue $Q_i(t)$}: Number of queued vehicles on approach lane $i$ with instantaneous velocity $v < 0.1\,\text{m/s}$:
    \begin{equation}
        Q_i(t) = \sum_{v \in \mathcal{V}_i(t)} \mathbb{I}(v_v(t) < 0.1\,\text{m/s})
    \end{equation}
    \item \textbf{Advance Upstream Inflow $A_i(t)$}: Vehicle count in the upstream entry segment ($x_v < 100.0\,\text{m}$), averaged over rolling window $W_A = 20\,\text{s}$ and scaled to equivalent arrival rate ($\text{veh/s}$):
    \begin{equation}
        A_i(t) = \frac{1}{10 \cdot W_A} \sum_{\tau=t-W_A}^t \sum_{v \in \mathcal{V}_i(\tau)} \mathbb{I}(x_v(\tau) < 100.0\,\text{m})
    \end{equation}
    \item \textbf{Queue Growth Rate $g_i(t)$}: Backward finite difference over window $W_G = 10\,\text{s}$:
    \begin{equation}
        g_i(t) = \frac{Q_i(t) - Q_i(t - W_G)}{W_G}
    \end{equation}
    \item \textbf{Accumulated Waiting Burden $W_i(t)$}: Average cumulative waiting duration across halting vehicles:
    \begin{equation}
        W_i(t) = \frac{\sum_{v \in \mathcal{V}_i(t), v_v < 0.1} w_v(t)}{\max(Q_i(t), 1)}
    \end{equation}
    \item \textbf{Continuous Red Age $R_d(t)$}: Elapsed continuous seconds that direction $d$ has been red while experiencing non-zero queue ($Q_d(t) > 0$).
\end{enumerate}

\section{Proposed Adaptive Signal Controller}
\subsection{Dimensionless Feature Normalization}
To ensure scale-invariance across heterogeneous traffic metrics, all state variables are normalized by physical engineering reference scales:
\begin{align}
    q_i(t) &= \min\left( \frac{Q_i(t)}{Q_{\text{ref}}}, 2.0 \right), \quad Q_{\text{ref}} = 20.0\,\text{veh} \\
    a_i(t) &= \min\left( \frac{A_i(t)}{A_{\text{ref}}}, 2.0 \right), \quad A_{\text{ref}} = 0.40\,\text{veh/s} \\
    w_i(t) &= \min\left( \frac{W_i(t)}{W_{\text{ref}}}, 2.0 \right), \quad W_{\text{ref}} = 30.0\,\text{s} \\
    \hat{g}_i(t) &= \max\left( -1.0, \min\left( 1.0, \frac{g_i(t)}{G_{\text{ref}}} \right) \right), \quad G_{\text{ref}} = 0.50\,\text{veh/s} \\
    r_d(t) &= \min\left( \frac{R_d(t)}{R_{\text{ref}}}, 1.0 \right), \quad R_{\text{ref}} = 50.0\,\text{s}
\end{align}
Directional queue growth combines only positive expansion: $g_d(t) = \sum_{i \in d} \max(0.0, \hat{g}_i(t))$.

\subsection{Composite Directional Priority Score}
The composite priority score $S_d(t)$ for direction $d \in \{\text{NS}, \text{EW}\}$ is computed as:
\begin{equation}
    S_d(t) = \alpha \sum_{i \in d} q_i(t) + \beta \sum_{i \in d} a_i(t) + \gamma g_d(t) + \delta \sum_{i \in d} w_i(t) + \eta r_d(t)
    \label{eq:priority_score}
\end{equation}
with frozen parameter weights established via independent calibration: $\alpha = 1.00, \beta = 0.75, \gamma = 0.10, \delta = 0.25, \eta = 0.25$.

\subsection{Constrained Decision Logic and Hysteresis Barrier}
Let $d_{\text{curr}}$ be the active green phase and $d_{\text{opp}}$ the conflicting red phase. At each step $t$, the controller evaluates four sequential decision gates:
\begin{enumerate}
    \item \textbf{Minimum Green Gate}: If $t_{\text{green}} < G_{\text{min}} = 15\,\text{s}$, green is strictly maintained.
    \item \textbf{Maximum Green Gate}: If $t_{\text{green}} \ge G_{\text{max}} = 50\,\text{s}$ and $Q_{d_{\text{opp}}}(t) > 0$, transition to yellow is enforced.
    \item \textbf{Starvation Override Gate}: If $R_{d_{\text{opp}}}(t) \ge R_{\text{starve}} = 45\,\text{s}$ and $Q_{d_{\text{opp}}}(t) > 0$, transition to yellow is enforced.
    \item \textbf{Hysteresis Switching Condition}: If none of Gates 1–3 trigger:
    \begin{equation}
        \text{Action}(t) = \begin{cases}
            \text{Initiate Yellow Clearance } (Y=3\,\text{s}) & \text{if } S_{d_{\text{opp}}}(t) > S_{d_{\text{curr}}}(t) + H \\
            \text{Maintain Current Green} & \text{otherwise}
        \end{cases}
    \end{equation}
    where $H = 0.20$ is the calibrated dimensionless hysteresis barrier.
\end{enumerate}

\section{Experimental Methodology}
\subsection{SUMO Simulation Environment and Network Setup}
Simulations were conducted in Eclipse SUMO v1.27.1 in headless mode. The orthogonal network has four incoming approaches ($N, S, E, W$) of $250\,\text{m}$ length, single-lane geometry, and Krauss car-following dynamics. Vehicle arrivals follow Poisson processes specified in immutable route files generated prior to evaluation.

\subsection{Experimental Design and Cohort Isolation}
The evaluation matrix incorporates 4 controllers $\times$ 7 demand scenarios $\times$ 5 evaluation seeds ($k \in \{42, 43, 44, 45, 46\}$), yielding exactly 140 simulation runs of $1000\,\text{s}$ duration each. All performance metrics are extracted strictly from post-warmup \textbf{Cohort A} vehicles ($t_{\text{depart}} \ge 200\,\text{s} \land t_{\text{arrive}} \ge 200\,\text{s}$).

\subsection{Benchmarked Signal Controllers}
\begin{itemize}
    \item \textbf{C0 (Original Fixed-Time)}: 90-second static cycle (NS Green 42s, Yellow 3s, EW Green 42s, Yellow 3s).
    \item \textbf{C1 (Webster-Calibrated Fixed-Time)}: Constrained 36-second static cycle (NS Green 15s, Yellow 3s, EW Green 15s, Yellow 3s), derived from Webster's optimum formula.
    \item \textbf{C2 (Basic Queue Adaptive)}: Naive real-time queue comparison ($Q_{\text{NS}}$ vs $Q_{\text{EW}}$) with $G_{\text{min}}=15\text{s}, G_{\text{max}}=50\text{s}, R_{\text{starve}}=45\text{s}$, but no arrival sensing, queue growth, waiting burden, red aging, or hysteresis ($H=0$).
    \item \textbf{C3 (Proposed Multi-Factor Adaptive)}: The complete proposed architecture.
\end{itemize}

\subsection{Inferential Statistical Protocol}
The experimental unit is the matched simulation realization $(s, k)$, establishing $N = 5$ paired observations per scenario. Hypotheses were evaluated using two-sided paired Student's $t$-tests ($df = 4$), exact two-sided Wilcoxon signed-rank tests, and paired Cohen's $d_z$. To control for family-wise Type I error inflation across the 21 confirmatory primary tests ($7 \text{ scenarios} \times 3 \text{ comparators}$ on waiting time), both Holm-Bonferroni (FWER) and Benjamini-Hochberg (FDR) corrections were applied at $\alpha = 0.05$.

\section{Results and Analysis}
\subsection{Primary Outcome: Cohort A Waiting Time}
The complete confirmatory results for Cohort A Mean Waiting Time across the 21 primary hypothesis tests are summarized in Table~\ref{tab:primary_waiting_time}.

\begin{table}[htbp]
\centering
\scriptsize
\caption{Confirmatory Statistical Evaluation of Cohort A Mean Waiting Time ($N=5$ matched realizations per scenario)}
\label{tab:primary_waiting_time}
\begin{tabularx}{\textwidth}{lccccccccX}
\toprule
\textbf{Scenario} & \textbf{Comp.} & \textbf{C3 Mean} & \textbf{Comp Mean} & \textbf{Mean Diff (s)} & \textbf{\% Change} & \textbf{95\% CI (s)} & \textbf{Raw $p$} & \textbf{BH $q$} & \textbf{Classification} \\
\midrule
\textbf{low} & C0 & 4.32 & 11.94 & -7.62 & -63.8\% & [-8.48, -6.76] & 1.59e-05 & 0.00011 & Supported improvement \\
\textbf{low} & C1 & 4.32 & 4.56 & -0.24 & -5.3\% & [-1.25, 0.77] & 0.54203 & 0.66957 & Unsupported improvement \\
\textbf{low} & C2 & 4.32 & 4.29 & +0.02 & +0.6\% & [-0.75, 0.80] & 0.93546 & 0.98223 & Unsupported degradation \\
\midrule
\textbf{medium} & C0 & 8.40 & 16.40 & -8.01 & -48.8\% & [-11.06, -4.95] & 0.00189 & 0.00795 & Supported improvement \\
\textbf{medium} & C1 & 8.40 & 6.64 & +1.75 & +26.4\% & [-0.60, 4.10] & 0.10697 & 0.20422 & Unsupported degradation \\
\textbf{medium} & C2 & 8.40 & 6.81 & +1.59 & +23.4\% & [-0.76, 3.94] & 0.13314 & 0.21508 & Unsupported degradation \\
\midrule
\textbf{high} & C0 & 49.66 & 75.62 & -25.96 & -34.3\% & [-39.48, -12.43] & 0.00597 & 0.02089 & Supported improvement \\
\textbf{high} & C1 & 49.66 & 43.22 & +6.44 & +14.9\% & [0.86, 12.02] & 0.03269 & 0.06865 & Unsupported degradation \\
\textbf{high} & C2 & 49.66 & 44.91 & +4.75 & +10.6\% & [-2.68, 12.18] & 0.15028 & 0.22542 & Unsupported degradation \\
\midrule
\textbf{oversaturated} & C0 & 71.18 & 86.61 & -15.43 & -17.8\% & [-28.73, -2.13] & 0.03225 & 0.06865 & Unsupported improvement \\
\textbf{oversaturated} & C1 & 71.18 & 67.87 & +3.31 & +4.9\% & [-1.40, 8.03] & 0.12295 & 0.21508 & Unsupported degradation \\
\textbf{oversaturated} & C2 & 71.18 & 71.18 & 0.00 & 0.0\% & [-5.75, 5.75] & 1.00000 & 1.00000 & No detectable difference \\
\midrule
\textbf{asym\_ns} & C0 & 13.38 & 59.55 & -46.17 & -77.5\% & [-50.36, -41.98] & 6.80e-06 & 0.00011 & Supported improvement \\
\textbf{asym\_ns} & C1 & 13.38 & 36.31 & -22.93 & -63.2\% & [-37.36, -8.50] & 0.01157 & 0.03472 & Supported improvement \\
\textbf{asym\_ns} & C2 & 13.38 & 16.17 & -2.80 & -17.3\% & [-7.47, 1.87] & 0.17157 & 0.24020 & Unsupported improvement \\
\midrule
\textbf{asym\_ew} & C0 & 12.51 & 67.51 & -55.01 & -81.5\% & [-60.68, -49.33] & 1.13e-05 & 0.00011 & Supported improvement \\
\textbf{asym\_ew} & C1 & 12.51 & 32.60 & -20.09 & -61.6\% & [-33.72, -6.47] & 0.01493 & 0.03919 & Supported improvement \\
\textbf{asym\_ew} & C2 & 12.51 & 18.17 & -5.67 & -31.2\% & [-15.70, 4.37] & 0.19214 & 0.25219 & Unsupported improvement \\
\midrule
\textbf{dynamic} & C0 & 30.22 & 58.54 & -28.32 & -48.4\% & [-36.56, -20.09] & 0.00067 & 0.00353 & Supported improvement \\
\textbf{dynamic} & C1 & 30.22 & 30.82 & -0.60 & -1.9\% & [-3.91, 2.72] & 0.64262 & 0.71027 & Unsupported improvement \\
\textbf{dynamic} & C2 & 30.22 & 31.21 & -0.99 & -3.2\% & [-5.69, 3.71] & 0.59063 & 0.68907 & Unsupported improvement \\
\bottomrule
\end{tabularx}
\end{table}

\subsection{Analysis by Demand Regime}
\subsubsection{Directional Asymmetry (asym\_ns, asym\_ew)}
The proposed controller achieves its most pronounced and statistically supported improvements in asymmetric traffic regimes. Compared to Webster-calibrated fixed-time C1, waiting time decreases by 63.2\% in \texttt{asym\_ns} ($q = 0.0347$) and by 61.6\% in \texttt{asym\_ew} ($q = 0.0392$). Because C1 enforces equal 15-second green allocations to both directions, heavy arterial flows experience severe cycle failures. C3 dynamically allocates green capacity to the dominant approach while honoring minimum green on the minor cross street. Against C2, waiting times decrease by 17.3\% and 31.2\%, though non-significant after FDR correction due to sample size constraints ($N=5$).

\subsubsection{Symmetric Balanced Regimes (low, medium, high)}
Under balanced steady-state demands, C3 significantly outperforms uncalibrated C0 ($-34.3\%$ to $-63.8\%$), but does \textbf{not} outperform Webster-calibrated C1. In \texttt{low}, C3 and C1 are statistically indistinguishable ($4.32\,\text{s}$ vs $4.56\,\text{s}, p = 0.542$). In \texttt{medium} and \texttt{high}, C1 achieves lower waiting times than C3 ($1.75\,\text{s}$ lower in \texttt{medium}, $6.44\,\text{s}$ lower in \texttt{high}), although neither difference achieves significance after multiple-testing correction ($q = 0.204$ and $q = 0.069$). Under balanced flows, C1's rapid 36-second cycle has zero detector latency, whereas C3's green extension logic occasionally holds green past queue dissipation.

\subsubsection{Dynamic Surge Demand}
Under time-varying surge demand, C3 strongly outperforms C0 ($-48.4\%, q = 0.0035$) and matches C1 ($30.22\,\text{s}$ vs $30.82\,\text{s}$) and C2 ($31.21\,\text{s}$). However, C3 achieves this waiting time with 17.3\% fewer phase switches than C1 and 7.9\% fewer phase switches than C2, demonstrating superior stability during surge onset and dissipation.

\subsubsection{Oversaturated Regime}
Under oversaturated demand ($p=0.18$), arrival flow exceeds physical intersection capacity. C3 reduces waiting time by 17.8\% relative to C0, though this does not survive multiple testing correction ($q = 0.0686$). Between C3 and C2, waiting times are identical ($71.18\,\text{s}$ vs $71.18\,\text{s}, \Delta = 0.00\,\text{s}, p = 1.000$). However, C3 executes 4.6\% fewer phase switches ($p = 0.0078$).

\subsection{Secondary Outcomes: Signal Phase Switching Stability}
Across all seven scenarios, C3 achieves a statistically significant reduction in phase switch transitions compared to both C1 and C2 ($p < 0.028$ everywhere):
\begin{itemize}
    \item \texttt{low}: $-17.60$ switches vs C2 ($-20.3\%, p = 0.00019, d_z = -5.93$)
    \item \texttt{medium}: $-11.60$ switches vs C2 ($-11.5\%, p = 0.00094, d_z = -3.91$)
    \item \texttt{high}: $-7.60$ switches vs C2 ($-7.1\%, p = 0.00611, d_z = -2.37$)
    \item \texttt{oversaturated}: $-4.80$ switches vs C2 ($-4.6\%, p = 0.00776, d_z = -2.21$)
    \item \texttt{asym\_ns}: $-13.20$ switches vs C2 ($-18.3\%, p = 0.02719, d_z = -1.52$)
    \item \texttt{asym\_ew}: $-12.00$ switches vs C2 ($-16.6\%, p = 0.01673, d_z = -1.77$)
    \item \texttt{dynamic}: $-7.80$ switches vs C2 ($-7.9\%, p = 0.00030, d_z = -5.26$)
\end{itemize}
These results provide rigorous empirical proof that the composite objective's queue growth smoothing and hysteresis terms successfully suppress phase hunting without compromising throughput.

\section{Discussion}
\subsection{Operational Trade-Offs Across Traffic Regimes}
The empirical findings refute the notion that any single signal control algorithm is globally optimal across all traffic states. Fixed-time control calibrated via Webster's formula provides an exceptionally strong baseline in balanced steady-state regimes, where cycle predictability and rapid cycling (36 s) minimize delay. However, fixed timing suffers acute failure under directional imbalance. The proposed adaptive controller addresses this vulnerability, delivering $> 60\%$ waiting time reductions in asymmetric conditions while maintaining switching stability.

\subsection{Mechanism of Hysteresis Stabilization}
In naive adaptive controllers (such as C2), signal state toggles whenever the opposing queue exceeds the active queue by a single vehicle. When demands are balanced, queues cross this zero-difference boundary repeatedly, inducing rapid phase alternation. Each transition incurs 3 seconds of yellow lost time ($L = 6\,\text{s}$ per cycle). The hysteresis threshold $H = 0.20$ enforces a stabilization buffer, requiring opposing pressure to convincingly exceed the active phase before terminating green. This saves between 4.8 and 17.6 clearance transitions per 1000 seconds.

\section{Practical Deployment Considerations}
While evaluated in SUMO simulation, the proposed formulation was explicitly designed for real-world deployment on edge compute platforms (e.g., standard NEMA TS2 or 2070 traffic signal controllers).
\begin{enumerate}
    \item \textbf{Sensing Modalities}: In a field installation, halting queues and arrival counts can be obtained via stop-bar video cameras, thermal imaging, radar sensors, or advanced loop detector arrays. Upstream arrivals ($x < 100\,\text{m}$) can be captured using advance detection loops or connected vehicle basic safety messages (BSMs).
    \item \textbf{Computational Latency}: The average computation time per decision step is under $0.5\,\text{ms}$, well within the 1000 ms control step budget of commercial traffic signal master controllers.
    \item \textbf{Fail-Safe Fallback}: In the event of sensor occlusion or detector failure, the controller can seamlessly fall back to the Webster-calibrated 36-second fixed-time timing plan (C1).
\end{enumerate}

\section{Limitations and Future Work}
This study is subject to several specific limitations:
\begin{enumerate}
    \item \textbf{Simulation-Only Evaluation}: All findings reflect microscopic simulation in Eclipse SUMO; physical hardware sensor noise, adverse weather camera degradation, and communication latency were not experimentally tested.
    \item \textbf{Single-Intersection Topology}: The evaluation focused on an isolated 4-arm intersection; corridor coordination, progression offsets, and multi-intersection shockwave spillbacks require future investigation.
    \item \textbf{Absence of Vulnerable Road Users}: Pedestrian actuations, bicycle crossings, and emergency vehicle preemption were not incorporated into the decision logic.
    \item \textbf{Sample Size Power ($N=5$)}: While matched pairing isolates treatment effects, a sample size of $N=5$ realizations limits statistical power to detect subtle differences ($|d_z| < 1.0$), and prevents two-sided Wilcoxon tests from achieving $p < 0.05$ due to the mathematical rank floor ($p_{\text{min}} = 0.0625$).
\end{enumerate}

\section{Conclusions}
This research formulated, audited, and empirically evaluated a multi-factor adaptive traffic signal controller incorporating queue state, arrival pressure, queue growth finite differences, waiting burden, red-age fairness, and switching hysteresis. Evaluated across 140 simulation runs on a blind test set in Eclipse SUMO, the findings confirm that the controller's primary benefit is concentrated under directional demand asymmetry (reducing waiting times by $61.6\%$ to $63.2\%$ relative to calibrated fixed timing) and in universal phase switching reduction ($4.6\%$ to $20.3\%$ fewer transitions across all scenarios). Under balanced steady-state demand, calibrated fixed timing remains highly competitive. Future work will extend the formulation to multi-intersection arterial networks and incorporate pedestrian clearance constraints.

\section*{Declarations}
\textbf{Funding}: This research received no external grant funding.\\
\textbf{Conflict of Interest}: The authors declare that they have no competing financial or personal interests.\\
\textbf{Data Availability}: All simulation network configurations, immutable demand routes, and raw per-seed output datasets are available in the public project repository.\\
\textbf{Author Contributions}: Conceptualization, methodology, software implementation, validation, formal analysis, and manuscript preparation were executed by the research team.

\bibliographystyle{IEEEtran}
\begin{thebibliography}{99}
\bibitem{papageorgiou2003review}
M.~Papageorgiou, C.~Diakaki, V.~Dinopoulou, A.~Kotsialos, and Y.~Wang, ``Review of road traffic control strategies,'' \emph{Proceedings of the IEEE}, vol.~91, no.~12, pp. 2043--2067, 2003. \url{https://doi.org/10.1109/JPROC.2003.819610}

\bibitem{webster1958traffic}
F.~V. Webster, ``Traffic signal settings,'' \emph{Road Research Technical Paper No. 39}, Road Research Laboratory, London, 1958.

\bibitem{hunt1981scoot}
P.~B. Hunt, D.~I. Robertson, R.~D. Bretherton, and R.~I. Winton, ``SCOOT - a traffic responsive method of coordinating traffic signals,'' \emph{TRRL Laboratory Report 1014}, Transport and Road Research Laboratory, Crowthorne, UK, 1981.

\bibitem{sims1980sydney}
A.~G. Sims and K.~W. Dobinson, ``The Sydney coordinated adaptive traffic (SCATS) system---Philosophy and benefits,'' \emph{IEEE Transactions on Vehicular Technology}, vol.~29, no.~2, pp. 130--137, 1980. \url{https://doi.org/10.1109/T-VT.1980.23843}

\bibitem{gartner1983opac}
N.~H. Gartner, ``OPAC: A strategy for online computer control of traffic signals,'' \emph{Transportation Research Record}, vol. 906, pp. 1--7, 1983.

\bibitem{mirchandani2001real}
P.~Mirchandani and L.~Head, ``A real-time traffic-responsive signal control system: RHODES,'' \emph{Transportation Research Part C: Emerging Technologies}, vol.~9, no.~6, pp. 415--432, 2001. \url{https://doi.org/10.1016/S0968-090X(00)00047-4}

\bibitem{varaiya2013max}
P.~Varaiya, ``Max pressure control of a network of signalized intersections,'' \emph{Transportation Research Part C: Emerging Technologies}, vol.~36, pp. 177--195, 2013. \url{https://doi.org/10.1016/j.trc.2013.08.014}

\bibitem{wongpiromsarn2012distributed}
T.~Wongpiromsarn, T.~Uthaicharoenpong, Y.~Wang, E.~Frazzoli, and D.~Wang, ``Distributed model predictive control for traffic signal networks,'' \emph{IEEE Transactions on Intelligent Transportation Systems}, vol.~13, no.~2, pp. 635--648, 2012. \url{https://doi.org/10.1109/TITS.2011.2180382}

\bibitem{kouvelas2014maximum}
A.~Kouvelas, M.~Saeedmanesh, and N.~Geroliminis, ``Maximum pressure traffic signal control using perimeter flow regulation in urban networks,'' \emph{Transportation Research Board 93rd Annual Meeting}, no. 14-4113, 2014.

\bibitem{cools2013self}
S.~B. Cools, C.~Gershenson, and S.~P. Hoogendoorn, ``Self-organizing traffic lights: A realistic simulation,'' in \emph{Self-Organizing Systems}, ser. Lecture Notes in Computer Science, vol. 8221, Springer, Berlin, Heidelberg, 2013, pp. 45--55. \url{https://doi.org/10.1007/978-3-642-54124-7_5}

\bibitem{gershenson2005self}
C.~Gershenson, ``Self-organizing traffic lights,'' \emph{Complex Systems}, vol.~16, no.~1, pp. 29--53, 2005.

\bibitem{arel2010reinforcement}
E.~Arel, C.~Liu, T.~Urbanik, and A.~G. Kohls, ``Reinforcement learning-based multi-agent system for network traffic signal control,'' \emph{IET Intelligent Transport Systems}, vol.~4, no.~2, pp. 128--135, 2010. \url{https://doi.org/10.1049/iet-its.2009.0070}

\bibitem{wei2018intellilight}
H.~Wei, G.~Zheng, H.~Yao, and Z.~Li, ``IntelliLight: A reinforcement learning approach for intelligent traffic light control,'' in \emph{Proceedings of the 24th ACM SIGKDD International Conference on Knowledge Discovery \& Data Mining}, 2018, pp. 2496--2505. \url{https://doi.org/10.1145/3219819.3220096}

\bibitem{wei2019survey}
H.~Wei, G.~Zheng, V.~Gayah, and Z.~Li, ``A survey on traffic signal control methods in the era of big data and AI,'' \emph{Frontiers in Big Data}, vol.~2, p.~38, 2019. \url{https://doi.org/10.3389/fdata.2019.00038}

\bibitem{genders2016using}
W.~Genders and S.~Razavi, ``Using deep reinforcement learning to model traffic signal control,'' \emph{Transportation Research Record}, vol. 2588, no.~1, pp. 26--34, 2016. \url{https://doi.org/10.3141/2588-04}

\bibitem{newell1989theory}
G.~F. Newell, ``Theory of highway traffic signals,'' \emph{Institute of Transportation Studies, University of California, Berkeley, Research Report UCB-ITS-RR-89-7}, 1989.

\bibitem{miller1963computer}
A.~J. Miller, ``A computer control system for traffic networks,'' in \emph{Proceedings of the Second International Symposium on the Theory of Road Traffic Flow}, London, 1963, pp. 200--220.

\bibitem{nema2005ts2}
National Electrical Manufacturers Association (NEMA), ``NEMA standards publication TS 2-2003: Traffic controller assemblies with NTCIP requirements,'' Rosslyn, VA, 2005.

\bibitem{hurlbert1984pseudoreplication}
S.~H. Hurlbert, ``Pseudoreplication and the design of ecological field experiments,'' \emph{Ecological Monographs}, vol.~54, no.~2, pp. 187--211, 1984. \url{https://doi.org/10.2307/1942661}

\bibitem{lopez2018microscopic}
P.~A. Lopez, M.~Behrisch, L.~Bieker-Walz, J.~Erdmann, Y.-P. Fl{\"o}tter{\"o}d, R.~Hilbrich, L.~L{\"u}cken, J.~Rummel, P.~Wagner, and E.~Wie{\ss}ner, ``Microscopic traffic simulation using SUMO,'' in \emph{21st International Conference on Intelligent Transportation Systems (ITSC)}, 2018, pp. 2575--2582. \url{https://doi.org/10.1109/ITSC.2018.8569938}

\bibitem{holm1979simple}
S.~Holm, ``A simple sequentially rejective multiple test procedure,'' \emph{Scandinavian Journal of Statistics}, vol.~6, no.~2, pp. 65--70, 1979.

\bibitem{benjamini1995controlling}
Y.~Benjamini and Y.~Hochberg, ``Controlling the false discovery rate: A practical and powerful approach to multiple testing,'' \emph{Journal of the Royal Statistical Society: Series B (Methodological)}, vol.~57, no.~1, pp. 289--300, 1995. \url{https://doi.org/10.1111/j.2517-6161.1995.tb02031.x}

\bibitem{cohen1988statistical}
J.~Cohen, \emph{Statistical Power Analysis for the Behavioral Sciences}, 2nd~ed.\hskip 1em plus 0.5em minus 0.4em\relax Hillsdale, NJ: Lawrence Erlbaum Associates, 1988.

\bibitem{trb2010hcm}
Transportation Research Board (TRB), \emph{Highway Capacity Manual (HCM 2010)}.\hskip 1em plus 0.5em minus 0.4em\relax Washington, D.C.: National Research Council, 2010.

\end{thebibliography}

\end{document}